\documentclass[11pt,reqno]{amsart}

\usepackage[utf8]{inputenc}
\usepackage{amsmath,amssymb,mathtools}
\usepackage{microtype}
\usepackage{xcolor}
\usepackage{hyperref}

\hypersetup{
  colorlinks=true,
  linkcolor=blue!55!black,
  citecolor=blue!55!black,
  urlcolor=blue!65!black
}

\newtheorem{theorem}{Theorem}[section]
\newtheorem{proposition}[theorem]{Proposition}
\newtheorem{lemma}[theorem]{Lemma}
\newtheorem{corollary}[theorem]{Corollary}
\theoremstyle{remark}
\newtheorem{remark}[theorem]{Remark}

\newcommand{\R}{\mathbb{R}}
\newcommand{\Qlx}{\mathcal{Q}}
\newcommand{\eps}{\varepsilon}
\newcommand{\dd}{\,dx}

\title[Intermediate $Q$-curvature positivity fails]
{An explicit counterexample to intermediate
$Q$-curvature positivity in dimension thirty}

\author{[Author Name]}
% \address{Department, University, Postal address}
% \email{author@example.edu}

\date{August 18, 2026}
\subjclass[2020]{53C18, 35J30, 35B09}
\keywords{conformal metric, $Q$-curvature, polyharmonic equation,
positivity, counterexample}

\begin{document}

\begin{abstract}
\textbf{Relation to the problem list.}
This manuscript addresses Question~50 (L50), ``Intermediate $Q$-curvature
positivity.''  It gives a full explicit counterexample in the stated class,
but it is not a new disproof: it is a self-contained specialization of the
V\'etois--Zeitler construction, which had already disproved the conjecture.

Li and Xu conjectured that, for a conformally Euclidean metric
\(g=u^{4/(n-2m)}|dx|^2\) on \(\R^n\), positivity and a slow-decay
barrier for the top-order curvature \(Q_g^{(2m)}\) force positivity of
every intermediate curvature \(Q_g^{(2k)}\), \(1\leq k<m\).
V\'etois and Zeitler recently disproved this statement for \(m\geq4\)
in sufficiently large dimensions.  We give a self-contained, completely
worked specialization of their construction to \((m,n)=(4,30)\).
For all sufficiently small \(\eps>0\), we construct a smooth positive
function \(u_\eps\) such that
\[
 g_\eps=u_\eps^{2/11}|dx|^2,
 \qquad 0<c_\eps\leq Q_{g_\eps}^{(8)}\leq C_\eps,
 \qquad Q_{g_\eps}^{(6)}(0)<0.
\]
The sign at the origin is certified by an exact sixth-order Taylor
calculation for a two-pole Riesz potential with weights \(10\) and \(1\).
All normalization factors and the slow-decay hypothesis are verified
explicitly.
\end{abstract}

\maketitle

\section{The conjecture and the counterexample}

We use the sign convention
\(\Delta=\sum_{j=1}^n\partial_j^2\), so that \(-\Delta\) is positive.
For a smooth positive function \(u\) on \(\R^n\), integers
\(2\leq m<n/2\) and \(1\leq k\leq m\), set
\begin{equation}\label{eq:LX-definition}
 \Qlx_{2k}[u]
 :=u^{-\frac{n+2k}{n-2m}}
   (-\Delta)^k\!\left(u^{\frac{n-2k}{n-2m}}\right).
\end{equation}
Thus the conformal metric
\[
 g=u^{\frac{4}{n-2m}}|dx|^2
\]
satisfies the equations used by Li and Xu \cite[(1.1)--(1.2)]{LiXu}:
\[
 (-\Delta)^k\!\left(u^{\frac{n-2k}{n-2m}}\right)
 =\Qlx_{2k}[u]\,u^{\frac{n+2k}{n-2m}}.
\]
Their Conjecture~1 asserts that if \(\Qlx_{2m}[u]>0\) and, for some
\(-2m<s\leq0\),
\begin{equation}\label{eq:barrier}
 \Qlx_{2m}[u](x)\geq c_0|x|^s
 \quad\text{for all sufficiently large }|x|,
\end{equation}
then \(\Qlx_{2k}[u]>0\) for every \(1\leq k<m\).

Our main result is the following explicit negation of that assertion.

\begin{theorem}\label{thm:main}
There is a smooth positive function \(u\) on \(\R^{30}\) for which
\[
 g=u^{2/11}|dx|^2
\]
has the following properties:
\begin{enumerate}
\item there are constants \(0<c<C<\infty\) such that
\begin{equation}\label{eq:top-bounded}
 c\leq \Qlx_8[u](x)\leq C
 \qquad\text{for every }x\in\R^{30};
\end{equation}
\item the intermediate curvature of order six is negative at the origin:
\begin{equation}\label{eq:q6-negative}
 \Qlx_6[u](0)<0.
\end{equation}
\end{enumerate}
In particular, \eqref{eq:barrier} holds with \(s=0\), and Li--Xu's
Conjecture~1 is false.
\end{theorem}

The geometric GJMS normalization used in
\cite[(1.1)]{VetoisZeitler} is
\begin{equation}\label{eq:standard-normalization}
 Q_g^{(2k)}=\frac{2}{n-2k}\,\Qlx_{2k}[u].
\end{equation}
The factors in \eqref{eq:standard-normalization} are positive.  Hence
Theorem~\ref{thm:main} is unchanged if all curvatures are interpreted in
that normalization.  In particular, when \((n,m)=(30,4)\),
\begin{align}
 Q_g^{(8)}
 &=\frac1{11}u^{-19/11}(-\Delta)^4u,
 \label{eq:q8-standard}\\
 Q_g^{(6)}
 &=\frac1{12}u^{-18/11}(-\Delta)^3(u^{12/11}).
 \label{eq:q6-standard}
\end{align}

The construction below is the \((m,n)=(4,30)\) instance of the
two-peak construction of V\'etois and Zeitler
\cite[Theorem~1.2 and \S4]{VetoisZeitler}.  The purpose of this note is
not to claim a new disproof, but to record one concrete family and give
an independent exact calculation of its negative sixth-order curvature.

\section{The two-bump conformal factor}

Let \(e_1=(1,0,\dots,0)\in\R^{30}\).  Choose
\(\rho\in C_c^\infty(B_1(0))\) such that
\begin{equation}\label{eq:mollifier}
 \rho\geq0,
 \qquad \int_{\R^{30}}\rho(x)\,dx=1,
 \qquad
 \rho_\eps(x):=\eps^{-30}\rho(x/\eps).
\end{equation}
Put
\begin{equation}\label{eq:h-definition}
 h_\eps(x)
 :=\eps(1+|x|^2)^{-19}
   +10\rho_\eps(x-e_1)+\rho_\eps(x+e_1),
 \qquad 0<\eps<\frac14.
\end{equation}
In particular, \(h_\eps\) is smooth and strictly positive everywhere.
Let \(\omega_{29}=|\mathbb S^{29}|\), and define
\begin{equation}\label{eq:kappa}
 \kappa
 :=\frac{\Gamma(12)}{2^7\,3!\,\Gamma(15)\omega_{29}}
 =\frac{11\Gamma(11)}{2^8\pi^{15}\Gamma(4)}
\end{equation}
and
\begin{equation}\label{eq:u-definition}
 u_\eps(x)
 :=\kappa\int_{\R^{30}}|x-y|^{-22}h_\eps(y)\,dy.
\end{equation}

\begin{lemma}\label{lem:fundamental}
For every \(0<\eps<1/4\), the function \(u_\eps\) is smooth and positive
on \(\R^{30}\), and
\begin{equation}\label{eq:top-equation}
 (-\Delta)^4u_\eps=11h_\eps
 \qquad\text{in }\R^{30}.
\end{equation}
\end{lemma}

\begin{proof}
The fundamental solution of \((-\Delta)^4\) on \(\R^{30}\) is
\begin{equation}\label{eq:fundamental-solution}
 \Phi(x)
 =\frac{\Gamma(11)}{2^8\pi^{15}\Gamma(4)}|x|^{-22}.
\end{equation}
Indeed, this is the specialization of
\[
 \frac{\Gamma((n-2m)/2)}{2^{2m}\pi^{n/2}\Gamma(m)}
 |x|^{2m-n}
\]
to \((n,m)=(30,4)\).  Equations
\eqref{eq:kappa}--\eqref{eq:u-definition} therefore give
\(u_\eps=11\Phi*h_\eps\), and \eqref{eq:top-equation} follows in the
sense of distributions.  The integral is finite and positive because
\(|x|^{-22}\) is locally integrable in dimension \(30\), while
\(h_\eps\) decays like \(|x|^{-38}\).  Local elliptic regularity applied
to \eqref{eq:top-equation} then gives \(u_\eps\in C^\infty(\R^{30})\).
\end{proof}

We next verify the global hypothesis in the conjecture.  Write
\begin{equation}\label{eq:W-psi}
 W(x):=(1+|x|^2)^{-11},
 \qquad
 \psi(x):=(1+|x|^2)^{-19}=W(x)^{19/11}.
\end{equation}

\begin{lemma}\label{lem:comparison}
For each fixed \(0<\eps<1/4\), there is a constant
\(C_\eps\geq1\) such that
\begin{equation}\label{eq:comparisons}
 C_\eps^{-1}\psi\leq h_\eps\leq C_\eps\psi,
 \qquad
 C_\eps^{-1}W\leq u_\eps\leq C_\eps W
 \quad\text{on }\R^{30}.
\end{equation}
\end{lemma}

\begin{proof}
The first comparison is immediate outside a fixed ball, where the two
mollified bumps vanish and \(h_\eps=\eps\psi\).  It extends over that ball
because \(h_\eps\) and \(\psi\) are smooth, \(h_\eps>0\), and
\(\psi>0\).

For completeness, the standard spherical bubble identity in the present
dimension is
\begin{equation}\label{eq:bubble-identity}
 (-\Delta)^4W
 =2^8\frac{\Gamma(19)}{\Gamma(11)}\,\psi.
\end{equation}
It follows either by four radial differentiations or from the conformal
covariance of \((-\Delta)^4\).  Convolving \eqref{eq:bubble-identity}
with \(\Phi\) from \eqref{eq:fundamental-solution} gives the Riesz-potential
identity
\begin{equation}\label{eq:background-potential}
 \kappa\int_{\R^{30}}|x-y|^{-22}\psi(y)\,dy
 =a_0W(x),
 \qquad
 a_0:=\frac{11\Gamma(11)}{2^8\Gamma(19)}>0.
\end{equation}
The possible polyharmonic difference vanishes by decay at infinity.

Thus the background term in \eqref{eq:u-definition} is
\(\eps a_0W\), which gives the lower bound for \(u_\eps\).  Each of the
two remaining potentials is smooth and bounded on compact sets.  Since
its source is compactly supported and has finite mass, it is bounded by
a constant times \(|x|^{-22}\) for large \(|x|\), hence by a constant
times \(W\) on all of \(\R^{30}\).  This proves the second comparison.
\end{proof}

\begin{proposition}\label{prop:top-curvature}
For each fixed \(0<\eps<1/4\), the metric
\begin{equation}\label{eq:g-epsilon}
 g_\eps:=u_\eps^{2/11}|dx|^2
\end{equation}
has strictly positive top-order curvature, uniformly bounded away from
zero and infinity:
\begin{equation}\label{eq:q8-bounds}
 0<c_\eps\leq \Qlx_8[u_\eps](x)\leq C_\eps<\infty
 \qquad(x\in\R^{30}).
\end{equation}
Consequently the slow-decay condition \eqref{eq:barrier} holds with
\(s=0\), and also with any fixed \(s\in(-8,0)\).
\end{proposition}

\begin{proof}
By \eqref{eq:LX-definition} and \eqref{eq:top-equation},
\begin{equation}\label{eq:q8-explicit}
 \Qlx_8[u_\eps]
 =11h_\eps u_\eps^{-19/11}>0.
\end{equation}
The two comparisons in \eqref{eq:comparisons}, together with
\(\psi=W^{19/11}\), give \eqref{eq:q8-bounds}.  If \(|x|\geq1\) and
\(s<0\), then \(|x|^s\leq1\), so the positive constant lower bound also
implies the barrier with rate \(s\).
\end{proof}

\section{The exact sixth-order sign calculation}

The limiting two-pole function is
\begin{equation}\label{eq:A-r}
 A_r(x):=r|x-e_1|^{-22}+|x+e_1|^{-22},
 \qquad r>0.
\end{equation}
Only its sixth-order Taylor polynomial at the origin is needed.

\begin{lemma}[Two-pole calculation]\label{lem:two-pole}
Let \(S=1+r\), \(\eta=(r-1)/(r+1)\), and
\begin{equation}\label{eq:F-r}
 F_r(x):=\left(\frac{A_r(x)}{S}\right)^{12/11}.
\end{equation}
Then
\begin{align}
 (-\Delta)^3F_r(0)
 ={}&125042688-460505088\eta^2
 \notag\\
 &\quad+557383680\eta^4-221921280\eta^6.
 \label{eq:eta-polynomial}
\end{align}
In particular, for \(r=10\),
\begin{equation}\label{eq:negative-exact}
 (-\Delta)^3F_{10}(0)
 =-\frac{43658772480}{1771561}<0.
\end{equation}
\end{lemma}

\begin{proof}
Write \(x=(t,z)\in\R\times\R^{29}\) and \(q=|z|^2\).  Assign weight
one to \(t\) and weight two to \(q\).  Since
\[
 |x\mp e_1|^{-22}=(1\mp2t+t^2+q)^{-11},
\]
the generalized binomial formula gives, through weighted degree six,
\begin{align}
 \frac{A_r}{S}
 ={}&1+22\eta t+(253t^2-11q)
      +\eta(2024t^3-264tq)\notag\\
 &+(12650t^4-3300t^2q+66q^2)\notag\\
 &+\eta(65780t^5-28600t^3q+1716tq^2)\notag\\
 &+(296010t^6-193050t^4q+23166t^2q^2-286q^3)
 +O_{\mathrm w}(7).
 \label{eq:A-expansion}
\end{align}
Here \(O_{\mathrm w}(7)\) contains only monomials of weighted degree at
least seven.  Expanding the \(12/11\)-power in \eqref{eq:F-r}, the part
of weighted degree six is
\begin{equation}\label{eq:F-degree-six}
 [F_r]_6=C_{60}t^6+C_{41}t^4q+C_{22}t^2q^2+C_{03}q^3,
\end{equation}
where direct collection of the generalized-binomial terms gives
\begin{align*}
 C_{60}&=396980+388056\eta^2-618240\eta^4+308224\eta^6,\\
 C_{41}&=-275460+17328\eta^2-26880\eta^4,\\
 C_{22}&=29148+2520\eta^2,\\
 C_{03}&=-364.
\end{align*}
These four coefficients can also be obtained by substituting
\eqref{eq:A-expansion} into
\[
 (1+X)^{12/11}
 =\sum_{j=0}^{6}\binom{12/11}{j}X^j+O_{\mathrm w}(7).
\]

For a function depending only on \((t,q)\), the Euclidean Laplacian on
\(\R\times\R^{29}\) is
\begin{equation}\label{eq:radial-laplacian}
 \Delta=\partial_t^2+58\partial_q+4q\partial_q^2.
\end{equation}
Only \([F_r]_6\) contributes to \(\Delta^3F_r(0)\), and
\begin{equation}\label{eq:laplacian-coefficients}
 \Delta^3F_r(0)
 =720C_{60}+4176C_{41}+43152C_{22}+1424016C_{03}.
\end{equation}
Substitution of the four displayed coefficients into
\eqref{eq:laplacian-coefficients}, followed by a change of sign, gives
\eqref{eq:eta-polynomial}.  Finally, when \(r=10\), one has
\(\eta=9/11\).  Substitution into \eqref{eq:eta-polynomial} yields
\[
 (-\Delta)^3F_{10}(0)
 =-\frac{43658772480}{11^6},
\]
and \(11^6=1771561\), proving \eqref{eq:negative-exact}.
\end{proof}

\begin{remark}[Connection with the palindromic sign polynomial]
\label{rem:K-polynomial}
In the notation arising from the sixfold integral calculation of
V\'etois--Zeitler \cite[\S4]{VetoisZeitler}, the specialization
\((m,n)=(4,30)\) gives
\[
 b_0=29952,
 \qquad b_1=-1009152,
 \qquad b_2=7168512
\]
and
\begin{equation}\label{eq:K-polynomial}
 K(r)=64r(b_0+b_1r+b_2r^2+b_1r^3+b_0r^4).
\end{equation}
The direct computation above is equivalently
\begin{equation}\label{eq:K-identity}
 (-\Delta)^3(A_r^{12/11})(0)
 =24(1+r)^{-54/11}K(r).
\end{equation}
For the integer choice \(r=10\),
\[
 b_0+10b_1+100b_2+1000b_1+10000b_0=-2842368,
\]
so \(K(10)=-1819115520<0\).  This provides an independent arithmetic
check on \eqref{eq:negative-exact}.
\end{remark}

We now transfer the point-mass calculation to the smooth metrics
\eqref{eq:g-epsilon}.

\begin{lemma}\label{lem:C6-convergence}
As \(\eps\downarrow0\),
\begin{equation}\label{eq:C6-convergence}
 u_\eps\longrightarrow\kappa A_{10}
 \quad\text{in }C^6(\overline{B_{1/4}(0)}).
\end{equation}
\end{lemma}

\begin{proof}
By \eqref{eq:background-potential}, the potential of the first term in
\eqref{eq:h-definition} is \(\eps a_0W\), which tends to zero in
\(C^\infty\).  For \(x\in\overline{B_{1/4}(0)}\) and
\(0<\eps<1/4\), the supports of
\(\rho_\eps(\,\cdot\mp e_1)\) stay a positive distance from \(x\).
The kernel \(|x-y|^{-22}\) and all of its \(x\)-derivatives are therefore
uniformly smooth on the relevant product sets.  The approximate-identity
property in \eqref{eq:mollifier} consequently gives convergence of the
two bump potentials, with every derivative, to
\(10|x-e_1|^{-22}\) and \(|x+e_1|^{-22}\), respectively.  This proves
\eqref{eq:C6-convergence}.
\end{proof}

\begin{proposition}\label{prop:q6-negative}
For all sufficiently small \(\eps>0\),
\begin{equation}\label{eq:q6-epsilon-negative}
 \Qlx_6[u_\eps](0)<0.
\end{equation}
More precisely,
\begin{equation}\label{eq:q6-limit}
 \lim_{\eps\downarrow0}\Qlx_6[u_\eps](0)
 =-\frac{43658772480}{1771561}
   (11\kappa)^{-6/11}<0.
\end{equation}
\end{proposition}

\begin{proof}
For \((n,m,k)=(30,4,3)\), definition \eqref{eq:LX-definition} reads
\begin{equation}\label{eq:q6-LX}
 \Qlx_6[u]
 =u^{-18/11}(-\Delta)^3(u^{12/11}).
\end{equation}
The limit \(\kappa A_{10}\) in \eqref{eq:C6-convergence} is bounded
away from zero on \(\overline{B_{1/4}(0)}\).  Composition with the smooth
map \(s\mapsto s^{12/11}\) therefore preserves the \(C^6\) convergence.
Since \(A_{10}(0)=11\) and
\(A_{10}^{12/11}=11^{12/11}F_{10}\), equations
\eqref{eq:q6-LX} and \eqref{eq:negative-exact} give
\begin{align*}
 \lim_{\eps\downarrow0}\Qlx_6[u_\eps](0)
 &=(11\kappa)^{-18/11}\kappa^{12/11}11^{12/11}
   (-\Delta)^3F_{10}(0)\\
 &=(11\kappa)^{-6/11}
   \left(-\frac{43658772480}{1771561}\right),
\end{align*}
which is \eqref{eq:q6-limit}.  Strict negativity persists for every
sufficiently small positive \(\eps\).
\end{proof}

\begin{proof}[Proof of Theorem~\ref{thm:main}]
Choose any \(\eps>0\) small enough for
Proposition~\ref{prop:q6-negative}, and set \(u=u_\eps\).
Lemma~\ref{lem:fundamental} gives smoothness and positivity.
Proposition~\ref{prop:top-curvature} gives
\eqref{eq:top-bounded} and the slow-decay barrier, while
Proposition~\ref{prop:q6-negative} gives \eqref{eq:q6-negative}.
\end{proof}

\section{Concluding remarks}

The counterexample has three features worth emphasizing.  First, the
top-order curvature is not merely positive: it is bounded between two
positive constants.  Second, the failure occurs for a smooth conformal
metric and is not an artifact of singular point masses; the point masses
enter only as a limit used to certify the sign.  Third, no numerical root
of the dimension-threshold polynomial is needed: the integer weight
\(r=10\) makes the sign strictly negative by exact arithmetic.

V\'etois and Zeitler prove a much broader result: for every \(m\geq4\),
they construct counterexamples in every sufficiently large dimension
\cite[Theorem~1.2]{VetoisZeitler}.  The calculation here isolates the
a particularly simple specialization, \((m,n)=(4,30)\), and keeps the
Li--Xu normalization visible throughout.

\begin{thebibliography}{99}

\bibitem{LiXu}
M.~Li and X.~Xu,
\emph{On positivity of the $Q$-curvatures of conformal metrics},
J. Funct. Anal. \textbf{289} (2025), no.~8, Paper No.~111011;
\href{https://doi.org/10.1016/j.jfa.2025.111011}
{doi:10.1016/j.jfa.2025.111011};
\href{https://arxiv.org/abs/2402.16277}{arXiv:2402.16277}.

\bibitem{VetoisZeitler}
J.~V\'etois and S.~Zeitler,
\emph{Positivity and non-positivity results for the sixth-order
$Q$-curvature of conformal metrics in $\mathbb{R}^n$},
preprint (20 July 2026),
\href{https://arxiv.org/abs/2607.18205}{arXiv:2607.18205v1}.

\end{thebibliography}

\end{document}
